\chapter*{Conferencia XVIII}
\addcontentsline{toc}{chapter}{LECTURA I - El timbre y el carácter del violín}
Estimados miembros del Club de Estudiantes de Violín: Esta hora marca el final de nuestro curso. Aquí tienen dos piezas de madera de interés: una es de pino blando de Michigan y la otra, de cedro blanco de Michigan. Ambas se presentan para demostrar que el aceite puede aplicarse incluso sobre madera blanda sin que penetre en ella. En los textos sobre violín, a menudo se afirma que el uso de aceite como agente protector es peligroso porque penetra en la madera. También se afirma que la penetración del barniz de Cremona en la madera es la razón de los valores tonales del violín de Cremona; además, que el barniz de Cremona es, disculpe, Sr. Promotor, un barniz de aceite. Debido a la información contradictoria en los textos sobre violín, no pude llegar a una conclusión sobre este punto tan importante sin realizar una prueba bajo mi propia observación.
Es evidente, sin necesidad de realizar ninguna prueba, que la penetración del aceite debe disminuir la rapidez con la que el resorte de madera regresa a su posición de reposo; por lo tanto, la penetración del aceite en la caja de resonancia debe disminuir la potencia sonora del violín. Razoné de la siguiente manera: si se puede aplicar aceite a estas dos muestras de madera blanda sin que penetren, entonces se puede aplicar aceite con seguridad a cualquier violín. Para que dicha prueba fuera fiable, utilicé aceite crudo, porque se seca más lentamente que el aceite hervido; además, mezclé aceite crudo con goma de mástique, la goma de secado más lento que conozco. No añadí nada más a esta mezcla. La superficie de la madera
274

Se alisó la superficie, pero no se empleó ningún relleno. La mezcla se aplicó con una almohadilla abrasiva, en capas finas. Se dejó secar cada capa antes de aplicar la siguiente.
Han pasado más de diez años desde que estas dos piezas de madera blanda recibieron este acabado. Con un cuchillo, raspo el acabado hasta llegar a la madera para demostrar que se aplicó suficiente mezcla para su protección. Con esta hoja afilada, retiro una fina viruta justo debajo del acabado. Como pueden observar, la madera está brillante y sin ninguna decoloración. Atribuyo este hecho a la forma de aplicación. He observado que una capa de aceite se seca con una rapidez proporcional a su grosor. Así, la capa más fina, posible con la almohadilla de frotar, se seca mucho más rápido que la capa más ligera posible con el pincel. Aplicado en capas finas, creo que la penetración en la madera tras la aplicación de barniz de aceite mediante frotamiento es menor que la penetración del alcohol tras la aplicación de barniz de alcohol con el pincel. Mis razones para preferir el barniz de aceite en el violín
son así:
1. Menor penetración en la madera.
2. Mayor atenuación.
3. Mayor elasticidad.
Lamentablemente, ahora hago una última crítica a la publicación. Esta cosa de apariencia inocente es siempre lista.
culo tanto para la pluma como para la lengua. Aunque
El puesto es a la vez objeto de burlas, mofas y críticas. 275

A pesar de las maldiciones y los insultos, su valor es incalculable. Su ajuste y reajuste ocupan los momentos de ocio de todo violinista del mundo. Aunque se mantiene en constante movimiento, la paciencia del mástil es infinita. Con partes anatómicas como la cabeza y el pie, el mástil se sostiene con igual firmeza en ambos extremos. En un lenguaje superlativo, los franceses superan a todos los demás pueblos al llamarlo "I'aime" (el alma).
¡Oh, tú, Poste! Tu mismo nombre
Nos haces pensar que eres dócil; sin embargo, eres precisamente lo que confiere el título al rey.
Pequeña eres, de poco espacio, y ningún valor aparece en tu rostro.
Sin embargo, eres un pensamiento precioso que jamás se compraría con las riquezas de la India. Llamarte "alma" no daña a nadie.
Es tuyo por derecho dar tono a A y E, que consideramos superior al valor del oro de Creso.
Ya hemos superado la época de la prohibición del alcohol. Como recordarán, se nos prometió algo para calmar la sed en este momento; y, si piensan como yo, celebraremos esta ocasión presentando nuestra fiesta de la Pascua judía.
¿Acaso no hemos superado dificultades inexpugnables?
¿No somos nosotros también "secos"?
En verdad, este polvo a lo largo del camino del violín es asfixiante.
No es mejor dedicar todo nuestro tiempo a la aridez;
276

Sin una sonrisa ocasional, la vida apenas merece la pena. Si el verdadero trabajo consiste en disfrutar del juego, entonces, con toda sinceridad, estamos preparados para un momento de recreación. En cuanto a la cantidad de trabajo que requiere, no conozco ninguna ocupación que supere el dominio de las posibilidades del violín. Incluso leer su partitura resulta agotador. Aunque se han dedicado 400 años a considerar esas posibilidades, esta cuestión sigue siendo nueva para cada nueva generación.
¿Seguirá siendo nuevo alguna vez? No hay duda.
El encanto del sonido del violín posee un poder hechizante que supera las cifras aritméticas. La ofrenda diaria de riquezas al "rey" supera todas las ofrendas a todos los demás hechiceros juntos. Si bien la teoría microbiana actual abarca la Tierra; si bien el cazador de gérmenes invade las fuentes mismas de la vida en busca de sus "cultivos"; si bien ha segregado y nombrado innumerables micrococos patógenos y saprófitos, sin embargo, de manera inexplicable, ha descuidado el violinicus universalis. Considerando la propensión epidémica de este germen, tal negligencia es asombrosa. Se encuentra en toda la Tierra y en todas las latitudes. No conoce prejuicios por credos, razas ni colores. Posiblemente, su negligencia por parte del microscopista se deba a que este germen es visible a la luz del sol, de la luna, de las estrellas, autoluminoso en la oscuridad y nunca sufre eclipses. El violinicus universalis es inerradicable. En todos los climas, su presencia se manifiesta mediante síntomas uniformes, a saber: intoxicación e indiferencia hacia la riqueza.
277

Ni la riqueza, ni el rango, ni la posición social garantizan la inmunidad ante este germen incontenible. El millonario y el mendigo se disputan un lugar bajo el arco del artista. Obsérvalos en el instante en que el primer tono dulce, tierno e intenso llega a sus oídos. Desde ese momento, ninguno se mueve, ninguno respira. Si tuvieran los ojos cerrados, permanecerían cerrados. Si tuvieran los labios entreabiertos, permanecerían entreabiertos. Ambos beben, beben el sonido más dulce de la tierra.
¿Intoxicado? Pregúntatelo.
Además, cuanto más beben, más quieren.
¡Música! ¡Música hermosa! En todo el mundo, desde la cabaña hasta el palacio, eres bienvenida. Tu lenguaje es comprendido por todos. Como el sol, tu presencia trae calidez. Tu poder para conmover el corazón humano no tiene igual. Tu elocuencia impone silencio, y el silencio es voluntario. Tus seguidores trascienden a la humanidad, llegando incluso al reino animal. Tu adoración es ilimitada e infinita, como el latido de los corazones.
¿Qué es la música? 
Para empezar, es una demostración de la facilidad con la que el idiota sonriente puede cuestionar al sabio. Pensando que quizás seas un sabio, y sin ver el peligro que se avecina, comienzas diciendo: "La música es algo construido a partir de la nada tangible; un producto del genio que trabaja en el éter para el abstemio; una entidad desconocida salvo para el oído; materia sin pesadez; material inconcebible a partir de lo intangible


reinos del pensamiento y y ." Como usted
Al ver la sonrisa cada vez más amplia de aquel idiota, apartas la mirada con incomodidad.
Según consta en los registros, la antigüedad de la música es inmensa. Estos registros transportan al estudiante casi hasta la época en que Adán, en el Edén, se convirtió en director de ópera.
Génesis iv:21. Hablando de los descendientes de Caín: Y el nombre de su hermano era Jubal. Él fue el padre de todos los que tocan el arpa y
órgano.
Así pues, afirmamos descender de Jubal. Nos enorgullece saber que nuestros antepasados dominaban el arpa, pues el arpa es de cuerdas.
[Solo echamos una mirada de reojo al órgano.] Nos deleitamos especialmente con aquellos descendientes de Jubal que tocan las cuerdas. Nos complacen sus voces suaves y su carácter tranquilo. Viven en un mundo aparte y dialogan con la mirada y las cuerdas. Es cierto que están lejanamente emparentados con Caín; también es cierto que «la sangre lo dirá»; sin embargo, jamás «asesinan la música», ni «provocan disturbios».
¿Podré ser perdonado alguna vez?
Que el crecimiento de la música, no hay nada en la historia que muestre igual deliberación. Es cierto que las cosas de crecimiento más lento duran más. Entre otros bienes, Salomón poseía una banda de cuatro mil trompetistas. Según la mejor luz disponible, las partituras interpretadas por esta magnífica banda se limitaban a cuatro tonos. El desarrollo de nuestras escalas mayores y menores requirió el tiempo desde Jubal hasta Palestrina en el siglo XVI de la era cristiana.
279

Era Ian. Le debemos mucho a Palestrina. Antes de Palestrina, esa maravillosa cosa llamada "armonía" no existía, ni se habría podido emplear aunque hubiera existido. En la construcción a gran escala, el genio de Palestrina hizo posible la orquesta moderna. Hoy, la pérdida de la orquesta sinfónica
sumiría al mundo en el luto. ¡Seis mil años después de Jubal!
En verdad, como ejemplo de crecimiento pausado, la música no tiene rival. Sin embargo, su desarrollo no es uniforme. Aún existen lugares donde el método de Jubal, con seis mil años de antigüedad, se conserva en su estado original. No conozco nada más interesante que el momento en que Jubal se enfrenta a una interpretación musical a cargo de la orquesta sinfónica actual.
¿Recuerdas la reciente presentación del teatro de los isleños del Pacífico Sur? Su orquesta de una sola pieza, su único tambor de tronco hueco, su único parche de tambor.
¿La baqueta, el único hijo de Jubal que manejaba la baqueta, su inimitable mirada de éxtasis?
Después de dominar sus da capos y dal segnos con mi muy dolcissimo le pregunté: "¿Podrías cambiar la tónica, por favor?"
No respondió ni una palabra.
No había necesidad.
Sus ojos, iluminados por el alma, me miraron fijamente, y de ellos emanaban señales miles de años más antiguas que las de Salomón. En esos orbes luminosos leí con claridad: «Esa es toda la clave que conozco».
Nuestro sentido musical es producto de seis mil años.
280

cultura del año. El sentido musical de este hijo de Jubal aún está en el punto cero. Para él, el monótono y lúgubre sonido de su primitivo instrumento de percusión (probablemente el órgano, Gén. iv:21) es suficiente para transportarlo a las regiones de los dulces sueños; y, desde su punto de vista onírico, el cambio a nuestras escalas mayores y menores, el cambio a la amplia gama de tonos melódicos, el cambio a 24 tonalidades, el cambio a las partes armónicas, el cambio a la interpretación mediante una partitura de diversos recursos, no es para él más que un cambio a un ruido incomprensible.
Para él, lo que está en el cielo es el infierno. Para nosotros, su "órgano" es el infierno.
Sin embargo, no se le puede negar el cielo.
Solo hay una petición alternativa: que las distancias en el cielo sean lo suficientemente grandes como para dar cabida a todos aquellos que puedan tocar el arpa y el órgano.
De este modo, todos podrán permanecer en la unión.
Pero aún hay otro asunto para reflexionar seriamente. Esas mismas distancias y planos por los que suplicamos pueden causarnos problemas. Así pues: mientras estamos aquí abajo, un rasgo de la humanidad es el egoísmo; y, cuando nos dejamos llevar por el egoísmo, seguramente intentaremos apoderarnos del puesto más alto. Como el egoísmo ciega a la introspección, la decepción aguarda a algunos de nosotros que pensamos que podemos manejar el arpa. El "maestro" puede sentirse seguro de entrar por la puerta más alta; pero, ¿se comparan los estándares terrenales con los celestiales? Esa es la cuestión. Incluso el arco del
"wa2sro" puede ser demasiado áspero para el oído de un ángel. 281	Eso

Intentaremos entrar por la puerta más alta posible. Es probable que a algunos nos rechacen. El maestro aquí solo puede entrar por una puerta inferior allá arriba. Podemos estar seguros de que el violín entrará por la puerta más alta, porque solo las cuerdas pueden acompañar la voz angelical. Pero tenemos motivos para la esperanza. Estos motivos señalan ciertas señales; y estas señales indican el camino. Así, el maestro al menos podrá llegar a la puerta más alta; que entre por ella depende de las precauciones.
Si el mundo le debe mucho a Palestrina, le debe aún más a la iglesia, que ni por un instante deja de fomentar la música bella.
Los valientes merecen vivir.
La lucha por la supervivencia que el constructor de violines moderno ha librado durante 200 años no tiene parangón en la historia. Si consideramos que solo dos rasgos humanos han sido la causa de esta prolongada lucha, nos asombra la magnitud y la profundidad de dichos rasgos. Específicamente, esos rasgos son:
Codicia por las ganancias.
Credulidad del consumidor.
En esta ocasión, la codicia por las ganancias es mucho menos relevante que la credulidad del consumidor; y, debido a que la misma codicia por las ganancias se muestra igualmente en otras líneas de comercio; pero, la misma credulidad no se encuentra en ningún otro consumidor. Ningún hecho es más evidente en la historia del violín que el hecho de que la credulidad humana es colosal. Desde que Stradivarius y Joseph Guarnerius señalaron el camino, innumerables violines, de igual valor tonal,
282

Se han construido; pero los constructores no han recibido ningún reconocimiento.
La razón por la que estos constructores no reciben el reconocimiento que merecen por violines con una calidad tonal comparable a la del mejor Stradivarius o al mejor "Joseph" es un asunto digno de reflexión. Por esta razón, solo hay tres categorías sobre las que se puede atribuir la responsabilidad, a saber:
1. El promotor de violines. 2. El consumidor de violines. 3. El constructor de violines.
El promotor de violines solo se interesa por los productos de aquellos constructores que han alcanzado la fama, porque de lo contrario, las ganancias son insignificantes. Para obtener grandes ganancias, el promotor depende únicamente del sentimiento de los consumidores; y el sentimiento de los consumidores es como la credulidad de los consumidores; y la credulidad de los consumidores es colosal. Luego vienen los constructores mismos, una procesión, una procesión fúnebre, aparentemente sí, en realidad.
¿Por qué tanto luto?
Señor constructor de violines, por la amabilidad que le tengo, mi bisturí está decidido a llegar al fondo de su problema. No se inmute. Es por su bien. Por lo tanto, diagnostico la causa de su aflicción.
ser esa enfermedad fatal llamada "imitación".
Fue la imitación lo que sepultó tus esperanzas.
JB Vuillaume encabeza la procesión. Cuando el violín Stradivarius se puso de moda, Vuillaume construyó violines Stradivarius; los vendió; los vendió a expertos en valores tonales; prueba suficiente de que el producto Vuillaume igualaba al producto Stradivarius.
¡Error fatal! 283

Si esos violines hubieran estado etiquetados honestamente como "JB Vuillaume, a Paris", entonces París hoy sería tan famosa como Cremona.
Desde Vuillaume, la procesión se alarga rápidamente; por consiguiente, se encuentran violines Stradivarius en todos los pueblos de Europa y América del Norte, y últimamente se oyen incluso en China.
Cuando los crédulos se ven engañados, a la credulidad le sigue la ira; y, cuando están enojados, los engañados no aceptarán tu imitación como un regalo, aunque su calidad tonal supere a la de los mejores de todos los tiempos.
Los valientes merecen vivir; y el día está amaneciendo rápidamente para aquellos valientes constructores de violines que se atreven a poner sus nombres en las etiquetas. Incluso en las altas esferas, se reconoce ahora el mérito del violín moderno, y la mano del promotor de violines antiguos está perdiendo valor. De una autoridad tan acreditada como Chas. Reade, de Londres, proviene la afirmación de que el mejor Stradivarius, sin barnizar, vale hoy tan solo 125 dólares. Esto es un reconocimiento de que existen violines con un valor tonal igual o incluso superior al del Stradivarius en abundancia. De hecho, los violines nuevos alcanzan precios inimaginables hace 200 años. El destino del "géiser" de Cremona es un rasgo triste que acompaña a este cambio de sentimiento. Como la palabra "géiser" se aplica a diferentes objetos, es necesario especificar que su uso en este contexto no se refiere a esos géiseres gigantescos del Parque Nacional de Yellowstone, sino al "géiser de Cremona". Hay una diferencia. El géiser de Yellowstone muestra la acción de aguas profundas. El "géiser de Cremona"


Representa la acción de aguas poco profundas. Si alguien alguna vez tuvo ocasión de rezar para liberarse de sus amigos, ese es Antonius Stradivarius. Un hombre es un hombre, nunca más, a veces menos.
Pronto el "cañón" de Cremona perderá su ocupación; y, por su destino, mi pluma derrama lágrimas de tinta negra. Sin nada en qué gastar su imaginación, el "cañón" de Cremona debe morir.
de apoplejía. Tristemente, nos apartamos de la repugnante exhibición de superficialidad del "espectador"; y nuestro grado de náuseas es exactamente proporcional a nuestro grado de estima por el objeto de su "espectáculo".
Un relato detallado del "chorro" de Cremona sería excesivo; una sobredosis inadmisible. Aunque dudo de la capacidad de mi pluma para ofrecer nuevas pistas sobre el "chorro" de Cremona, debido a su afán, se le permitirá intentarlo.
En tal juicio, es necesario citar un solo ejemplo de "efusión". Este ejemplo dice: "Él (Stradivarius) se casó con la rica viuda, Signora Capra, y después siguió su vocación elegida".
desde un punto de vista puramente artístico. ¡Por supuesto!
Lo mismo le ocurre a cualquier pobre luthier cuando se ve abatido por las exigencias de una fiebre recurrente. Aunque no fuera la intención del "goteo", esta afirmación prácticamente da a entender que Stradivarius conocía las exigencias de una fiebre recurrente antes de ser adoptado por aquella benevolente viuda.
Mi pluma señala el hecho de que el "chorro" ha ignorado por completo la equidad de la Signora Capra en
285

La gloria de Strad.
¡Joseph Guarnerius en la cárcel!
¡Y en la cárcel por una deuda insignificante! Es triste leer esto hoy en día, considerando que el constructor de violines moderno apenas presta atención a los \$500 por un solo violín; pero, el descuido del "acre" para la hija del carcelero supera
mera tristeza. Parece que incluso la Cremona
El "gusher" podría entender que la "divinidad que moldea nuestros fines", a falta de una viuda benevolente, puso a José en la cárcel con el propósito de brindarle los servicios de la hija del carcelero para conseguir el mejor material. Como "regular", estoy dispuesto a prescribirles a ustedes, pobres muchachos, lo siguiente: como tratamiento confiable y antifebril para la fiebre de "regreso rápido", tomen una viuda rica a la Stradivarius. Sinceramente espero que el "gusher" de Cremona pueda recuperar la oportunidad perdida de hacer justicia a los silenciosos socios de los genios de Cremona; pero, por el bien de los vivos, ruego sinceramente que el "gusher" le dé descanso a Stradivarius, el hombre honesto, concienzudo, ambicioso, incansable, amante del violín, violinista, fabricante de violines, nada más. 1644
1737 equivalía a todos los días del año mil novecientos cinco.
Se dice que el violín superior es producto del genio. El Honorable W. E. Gladstone afirmó que la evolución del violín requirió más reflexión que la evolución de la máquina de vapor.
En todas las actividades humanas, no conozco ninguna que posea igual dificultad para lograr la superioridad de
286

Producto con el arte de la construcción de violines. No conozco otro camino hacia la superioridad en la producción de violines que no sea la experiencia combinada con una aplicación intensa. Para mí, el violín superior aparece como producto de una habilidad mecánica superior combinada con un sentido musical superior y una aplicación intensa, todo dirigido a un material superior, nada más.
menos.
Que el arte de la construcción de violines, no conozco ninguna actividad humana que presente la misma tentación para el fraude. Dentro del violín reside gran parte del trabajo del que depende su valor; y, dado que gran parte de su interior permanece oculto, el fraude se hace presente. Nadie puede mirar solo el exterior de un violín y determinar que no hay fraude en su interior; sin embargo, todo comprador cree que el violín siempre mejora con el tiempo y el uso. Es la creencia universal de que el violín siempre mejora con el tiempo y el uso. Esta creencia universal en el engaño me duele el corazón. Tal creencia es una llama insensata que conduce diariamente a sus víctimas al abismo de la desesperación; y la procesión hacia allí es multitudinaria. En todo el mundo, con un volumen cada vez mayor, los sentimientos más tiernos de los corazones humanos se derraman a los pies de la música hermosa. Que se permita que el fraude prospere a costa de tales sentimientos enloquece al diablo de alegría.
Ustedes, que han leído atentamente estas páginas, están preparados para comprender las dificultades que rodean el camino del luthier concienzudo. Están preparados para el hecho de que el artesano más hábil, por mucho que lo intente, ocasionalmente debe
287

Se enfrentan a la derrota, en mayor o menor medida. La historia nos enseña que, incluso para el mejor artesano, el reconocimiento tarda en llegar. Ante esta realidad, ante las dificultades económicas y, desde mi punto de vista, aquel luthier que resiste con firmeza la tentación de defraudar es un héroe.
De esa materia están hechos los héroes.
Al recordar el pasado, al honrar la justicia y al amar la música, dondequiera y cuandoquiera que encuentres a tal héroe, te apresurarás a coronarlo con el honor mientras aún viva; y al otorgarle honores, al honrar la equidad, no olvidarás la justicia que le corresponde a su socio silencioso. Durante la penumbra de la espera, su voz ha sido su aliento. Concédele los honores que merece.
¡Hermosa música! Tú puedes salvar Cuando otros fallan. A mí me diste El poder, la voluntad, el pensamiento para mover De poner el surco del pecador obstinado. ¡Hermosa música! Tú eres la hoja Abriendo de par en par la puerta al cielo. ¿Te amo? Hasta el final,
Contigo al cielo, mi camino recorro. ¡Tú, Música! A tus pies yace mi tributo.
